%!TEX root = ../main.tex
% APÉNDICE C — Detalles técnicos e implementación

\chapter{Detalles técnicos e implementación}\label{apx:tecnicos}

\begin{guia}
Todo lo que alguien necesita para reproducir los resultados:
hiperparámetros exactos, identificadores y versiones de los datos,
versiones de librerías, variables de entorno (sin secretos) y esquemas
de base de datos. Reproducibilidad total es un marcador de excelencia.
\end{guia}

\section{Configuración de los modelos}

\begin{table}[H]
    \centering
    \caption{Hiperparámetros del modelo final (ejemplo de estructura)}
    \label{tab:hiperparametros}
    \small
    \begin{tabular}{l l l}
        \toprule
        \tb{Parámetro} & \tb{Valor} & \tb{Criterio de selección} \\
        \midrule
        \completar{n\_estimators} & \completar{500} & \completar{búsqueda en grilla, validación cruzada 5 pliegues} \\
        \completar{max\_depth} & \completar{8} & \completar{...} \\
        \bottomrule
    \end{tabular}
\end{table}

\section{Datos utilizados}

\completar{Identificador, fuente, versión, fecha de descarga, licencia y
tamaño de cada conjunto de datos.}

\section{Arquitectura de la aplicación}

\subsection{Stack tecnológico}

\begin{lstlisting}[caption={Dependencias principales (ejemplo)}, label={lst:dependencias}]
python == 3.12
scikit-learn == 1.5.0
fastapi == 0.111.0
postgresql == 16
\end{lstlisting}

\subsection{Esquema de base de datos}

\begin{lstlisting}[language=SQL, caption={Definición de tablas (ejemplo de código SQL)}, label={lst:esquema}]
CREATE TABLE usuario (
    id          SERIAL PRIMARY KEY,
    correo      TEXT NOT NULL UNIQUE,
    creado_en   TIMESTAMPTZ NOT NULL DEFAULT now()
);
\end{lstlisting}

\subsection{Variables de entorno}

\completar{Nombre, propósito y ejemplo de valor de cada variable. Nunca
incluir secretos reales.}
